%=======================================================================
% Sina Ebrahimi
% Curriculum Vitae (public copy for sinaebrahimi.github.io)
% Hard limit: three sides of A4. Keep it that way.
% Requires resume.cls in the same folder.
%=======================================================================

\documentclass{resume}

\usepackage[a4paper,left=0.55in,right=0.55in,top=0.5in,bottom=0.5in]{geometry}
\usepackage[utf8]{inputenc}
\usepackage[T1]{fontenc}
\usepackage[english]{babel}
\usepackage{enumitem}
\usepackage[hidelinks]{hyperref}
\usepackage{url}
\urlstyle{same}
\def\UrlBreaks{\do\/\do-\do.\do?\do=\do_\do\&}
\newcommand{\weblink}[2]{\href{#1}{\nolinkurl{#2}}}
\usepackage[none]{hyphenat}
\usepackage{ragged2e}
\justifying
\sloppy
\setlength{\emergencystretch}{3em}

% Tighter section rules than the class default, to win vertical space
\def\sectionskip{\smallskip}
\def\sectionlineskip{\smallskip}
\def\nameskip{\medskip}
\def\addressskip{\smallskip}

\newcommand{\entry}[2]{{\bf #1} \hfill {#2}\\}
\newcommand{\sub}[1]{{\em #1}\par}

\setlist[itemize]{leftmargin=1.1em, topsep=1pt, itemsep=0pt, parsep=0pt}

\name{Sina Ebrahimi}
\address{\small PhD in End-to-End Network Slicing for 5G/6G Networks \addressSep\ Coventry, United Kingdom \addressSep\ +44 7377 436 435}
\address{\small \href{mailto:ebrahimi.sina.1372@gmail.com}{ebrahimi.sina.1372@gmail.com} \addressSep\ \href{https://sinaebrahimi.github.io/}{sinaebrahimi.github.io} \addressSep\ \href{https://orcid.org/0000-0001-7349-7351}{ORCID 0000-0001-7349-7351} \addressSep\ \href{https://scholar.google.com/citations?user=Y9kFAooAAAAJ}{Google Scholar}}

\begin{document}

%-----------------------------------------------------------------------
\begin{rSection}{Research Profile}
Researcher in artificial intelligence for the management of 5G and 6G mobile networks, with eight years spanning operator core-network strategy, router and transport R\&D, connected-vehicle integration and doctoral research. My work develops multi-agent deep reinforcement learning methods that coordinate radio access, transport and core network resources as a single end-to-end system, so that network slices can meet service guarantees at the timescales real networks operate on. I am lead author of the reference survey for this subfield in \textit{IEEE Communications Surveys and Tutorials}, and I review for six journals and for conference tracks and workshops.
\end{rSection}

%-----------------------------------------------------------------------
\begin{rSection}{Education}

\entry{PhD, Electrical and Computer Engineering}{October 2022 to September 2026}
\sub{Centre for Future Transport and Cities, Coventry University, Coventry, United Kingdom}
\begin{itemize}
  \item Thesis: \textit{Satisfaction-Driven Hierarchical Multi-Agent Deep Reinforcement Learning for End-to-End Network Slicing}. Awarded September 2026.
  %5G/6G End-to-End Network Slicing Framework for Vertical Industries
  \item Supervisors: Dr Olivier C. L. Haas (Director of Studies), Prof. Soufiene Djahel, Dr Hui Zhou, Dr Faouzi Bouali.
  \item Held a fully funded doctoral studentship covering all tuition fees and an annual maintenance stipend.
\end{itemize}
\smallskip

\entry{MSc, Software Engineering (Computer Networks)}{September 2017 to January 2020}
\sub{Department of Electrical and Computer Engineering, Tarbiat Modares University, Tehran, Iran. GPA 3.54/4.0}
\begin{itemize}
  \item Thesis: \textit{A Resource Allocation Method for Network Slicing in Software-Defined 5G Networks}.
  \item Supervisors: Dr Behzad Akbari and Dr Nader Mokari, Wireless Innovation Laboratory.
  \item Awarded full state-funded tuition through Iran's national graduate entrance examination.
\end{itemize}
\smallskip

\entry{BSc, Software Engineering}{January 2012 to November 2016}
\sub{Department of Electrical and Computer Engineering, University of Zanjan, Zanjan, Iran. GPA 3.31/4.0}
\begin{itemize}
  \item Awarded full state-funded tuition through Iran's national undergraduate entrance examination.
\end{itemize}

\end{rSection}

%-----------------------------------------------------------------------
\begin{rSection}{Awards and Honours}
\begin{itemize}
  \item \textbf{Third Best Paper Award}, ITU Kaleidoscope 2020, the International Telecommunication Union's academic conference, for \textit{Digital Transformation via 5G: Deployment Plans} (2020).
  \item \textbf{Fully funded doctoral studentship}, Centre for Future Transport and Cities, Coventry University (2022 to 2026). Competitive award covering full tuition and an annual stipend.
  \item \textbf{State-funded university places by national competitive examination}, Iran. Full tuition awarded on examination performance for both the BSc at the University of Zanjan and the MSc at Tarbiat Modares University.
  \item \textbf{Formal commendation from the Technical Deputy}, Mobile Communication Company of Iran (MCI), for research on 5G technologies and standardisation supporting the national 5G migration programme (2019).
  \item \textbf{Travel awards}: UNICO I+D 6G programme, Government of Spain, for participation in the i2CAT Advanced Internet Technologies Winter PhD School, Barcelona (2024); University of Bristol (JOINER Project), for the ENSURE poster presentation at the 1st Seminar of Towards Edge Intelligence (TEI), Bristol (May 2026).
\end{itemize}
\end{rSection}

%-----------------------------------------------------------------------
\begin{rSection}{Publications}

\textbf{Journal articles}
\begin{enumerate}[leftmargin=1.3em, topsep=2pt, itemsep=2pt, parsep=0pt]
  \item \textbf{S. Ebrahimi}, H. Zhou, S. Djahel, F. Bouali and O. C. L. Haas, ``ENSURE: Multi-Timescale End-to-End Slice Resource Management for 6G Networks,'' \textit{IEEE Transactions on Cognitive Communications and Networking}, manuscript TCCN-TP-26-0585, major revision under review, 2026. Publisher: IEEE. No public record until acceptance; manuscript available on request.
  \item \textbf{S. Ebrahimi}, F. Bouali and O. C. L. Haas, ``Resource Management From Single-Domain 5G to End-to-End 6G Network Slicing: A Survey,'' \textit{IEEE Communications Surveys and Tutorials}, vol. 26, no. 4, pp. 2836--2866, 2024. Publisher: IEEE. \weblink{https://doi.org/10.1109/COMST.2024.3390613}{doi.org/10.1109/COMST.2024.3390613}
  \item A. Akhtarshenas, M. A. Vahedifar, N. Ayoobi, B. Maham, T. Alizadeh, \textbf{S. Ebrahimi} and D. L\'{o}pez-P\'{e}rez, ``Federated Learning: A Cutting-Edge Survey of the Latest Advancements and Applications,'' \textit{Computer Communications}, vol. 228, art. 107964, 2024. Publisher: Elsevier. \weblink{https://doi.org/10.1016/j.comcom.2024.107964}{doi.org/10.1016/j.comcom.2024.107964}
  \item M. Tajallifar, \textbf{S. Ebrahimi}, M. R. Javan, N. Mokari and L. Chiaraviglio, ``Energy-Efficient Task Offloading Under E2E Latency Constraints,'' \textit{IEEE Transactions on Communications}, vol. 70, no. 3, pp. 1711--1725, 2022. Publisher: IEEE. \weblink{https://doi.org/10.1109/TCOMM.2021.3132909}{doi.org/10.1109/TCOMM.2021.3132909}
\end{enumerate}

\smallskip
\textbf{Peer-reviewed conference papers}
\begin{enumerate}[leftmargin=1.3em, topsep=2pt, itemsep=2pt, parsep=0pt]
  \item \textbf{S. Ebrahimi}, F. Bouali and O. C. L. Haas, ``MARSRA: Mobility-Aware RAN Slicing Resource Allocation for Open RAN Deployments,'' \textit{IEEE International Workshop on Computer Aided Modeling and Design of Communication Links and Networks (CAMAD)}, Athens, Greece, 2024. Publisher: IEEE. \weblink{https://doi.org/10.1109/CAMAD62243.2024.10942928}{10.1109/CAMAD62243.2024.10942928}
  \item A. Zakeri, N. Gholipoor, M. Tajallifar, \textbf{S. Ebrahimi}, N. Mokari, M. R. Javan and A. R. Sharafat, ``Digital Transformation via 5G: Deployment Plans,'' \textit{ITU Kaleidoscope Conference}, Ha Noi, Vietnam, 2020. Third Best Paper Award. Publisher: IEEE and ITU. \weblink{https://doi.org/10.23919/ITUK50268.2020.9303177}{doi.org/10.23919/ITUK50268.2020.9303177}
  \item \textbf{S. Ebrahimi}, A. Zakeri, N. Mokari and B. Akbari, ``Joint Resource and Admission Management for Slice-enabled Networks,'' \textit{IEEE/IFIP Network Operations and Management Symposium (NOMS)}, Budapest, Hungary, 2020. Publisher: IEEE. \weblink{https://doi.org/10.1109/NOMS47738.2020.9110322}{doi.org/10.1109/NOMS47738.2020.9110322}
\end{enumerate}

\smallskip
\textbf{Book chapter}
\begin{enumerate}[leftmargin=1.3em, topsep=2pt, itemsep=2pt, parsep=0pt]
  \item \textbf{S. Ebrahimi}, A. Nouruzi, A. Akhtarshenas, M. Mosahebfard and S. Djahel, ``AI-based Resilience and Failure Recovery in 6G Networks,'' in \textit{AI-Enabled Wireless Communication: Artificial Intelligence and 6G Networks}, Elsevier and Morgan Kaufmann, invited chapter, under development, to be submitted late October 2026.
\end{enumerate}

\end{rSection}

%-----------------------------------------------------------------------
\begin{rSection}{Research and Professional Experience}

\entry{Doctoral Researcher and Research Assistant}{October 2022 to September 2026}
\sub{Centre for Future Transport and Cities, Coventry University, United Kingdom}
\begin{itemize}
  \item Designed ENSURE, a cooperative multi-agent deep reinforcement learning framework in which four agents manage radio access, transport, core and end-to-end slice resources on three distinct control timescales, using soft actor-critic and masked proximal policy optimisation.
  \item Introduced an action-masking formulation that makes infeasible slice allocations unreachable by construction rather than penalising them after the fact, allowing learned policies to respect service-level guarantees.
  \item Produced the field's first end-to-end survey of slicing resource management across all three network domains, together with an open-source literature analysis pipeline enabling reproduction and extension by other groups.
  % \item Built and released the full simulation codebase for both frameworks under open licence.
\end{itemize}

\entry{Research Assistant, V2X Systems Integration (Project CERTUS)}{January 2025 to July 2025}
\sub{Centre for Connected and Autonomous Automotive Research, HORIBA MIRA, Nuneaton, United Kingdom}
\begin{itemize}
  \item Integrated and validated Cohda Wireless MK6 cellular vehicle-to-everything units on StreetDrone connected and autonomous vehicle platforms for this Innovate UK funded project.
  \item Tested OBU-to-OBU C-V2X sidelink (PC5) and DSRC (ITS-G5), plus cellular 5G NR uplink/downlink for OBU-to-RSU connectivity; diagnosed failures from packet captures.
  \item Analysed steering telemetry in Python to estimate control-loop latency and assess deadband mitigation strategies.
\end{itemize}

\entry{Network Solutions Specialist, Router R\&D}{December 2021 to October 2022}
\sub{SINA Communication Systems, Tehran, Iran}
\begin{itemize}
  \item Tested carrier-grade L2/L3 switch/router NOS modules: Wireshark/tcpdump debugging and verification, GNS3 emulation, TRex traffic generation and VIAVI test devices; documented test procedures and results.
  \item Worked across OF-DPA, OpenFlow, FRRouting, ONOS, and Stratum for programmable transport hardware.
\end{itemize}

\entry{Network Performance Tools Engineer}{March 2020 to August 2021}
\sub{Network Quality and Supervision Department, RighTel Telecommunications, Tehran, Iran}
\begin{itemize}
  \item Built the operator's KPI extraction and visualisation pipeline over operations support system telemetry, automating daily network health reporting and downtime alerting.
  \item Deployed and maintained the supporting database and version control infrastructure.
\end{itemize}

\entry{5G Core Network Researcher}{January 2018 to July 2019}
\sub{Mobile Network Technology Solution and Strategy Department, Mobile Communication Company of Iran (MCI)}
\begin{itemize}
  \item Contributed to the national operator's 5G core architecture and migration study, covering the service-based architecture, network functions, slicing, SDN and NFV, delivered jointly with Tarbiat Modares University.
  \item Authored technical reports and delivered internal training on 3GPP standardisation and deployment/migration options to engineering and strategy teams.
\end{itemize}

\entry{Research and Teaching Assistant}{January 2018 to January 2020}
\sub{Wireless Innovation Laboratory, Tarbiat Modares University, Tehran, Iran}
\begin{itemize}
  \item Published three peer-reviewed papers and supported MSc researchers on network resource allocation projects.
  \item Teaching assistant for Advanced Computer Networks at MSc level.
\end{itemize}

\end{rSection}

%-----------------------------------------------------------------------
\begin{rSection}{Peer Review and Technical Programme Committee Service}
\textbf{Journal reviewer:} IEEE Transactions on Mobile Computing; IEEE Transactions on Cognitive Communications and Networking; IEEE Transactions on Network and Service Management; IEEE Open Journal of the Communications Society; IEEE Networking Letters; Nature Scientific Reports.\\[2pt]
\textbf{Conference paper reviewer (tracks, symposia and/or workshops):} IEEE INFOCOM 2025; IEEE ICC 2025 and 2026; IEEE WCNC 2025; IEEE NetSoft 2026; IEEE VTC2024-Fall and VTC2026-Fall; IEEE MedComNet 2025; IEEE FINE 2026.\\[2pt]
\textbf{Technical programme committee (TPC) member:} IEEE ICC 2025, \textit{Mobile and Wireless Networks Symposium}; IEEE ICC 2026, \textit{Open RAN workshop WS-01}; IEEE WCNC 2025, \textit{Open RAN workshop WS05}; IEEE VTC2024-Fall workshops: \href{https://sites.google.com/view/vtc2024-fall-ccavs/workshop-committee}{\textit{B5G/6G support for space/air/ground/marine/submarine cooperative, connected, and autonomous vehicles}} and \href{https://sites.google.com/view/vtc2024-fall-iiots/workshop-committee}{\textit{B5G/6G Support for Next-generation Industrial IoT}}; IEEE FINE 2026; MEET 2026.\\[2pt]
\textbf{Professional membership:} IEEE member since 2020, including the Communications Society, Vehicular Technology Society, Computer Society and the UK and Ireland Section; student member of The Institute of Telecommunications Professionals (ITP) since September 2026.
\end{rSection}

%-----------------------------------------------------------------------
\begin{rSection}{Invited Talks, Presentations and Outreach}
\begin{itemize}
  \item \textbf{\href{https://extend-lab.github.io/TEI/\#posters}{1st Seminar of Towards Edge Intelligence (TEI)}}, Bristol Digital Futures Institute (BDFI), University of Bristol, Bristol, UK, May 2026. Presented the poster ``\textit{ENSURE: Multi-Timescale E2E Slice Resource Management for 6G Networks},'' with travel awarded by the University of Bristol through the JOINER Project.
  \item \textbf{Coventry Centre for Doctoral Training in AI Community Event}, Coventry, April 2026. Poster on multi-timescale end-to-end slice resource management using multi-agent deep reinforcement learning.
  \item \textbf{Advanced Internet Technologies Winter PhD School, i2CAT}, Barcelona, November 2024. Presented deep reinforcement learning for Open RAN slicing, with travel supported by the Spanish UNICO I+D 6G programme.
  \item \textbf{Network Softwarization and Security Laboratory, University College Dublin}, September 2024. Invited seminar on the application of deep reinforcement learning to radio access network slicing.
  \item \textbf{IEEE CAMAD 2024}, Athens, October 2024, and \textbf{IEEE/IFIP NOMS 2020}, Budapest, April 2020. Conference paper presentations.
  \item \textbf{Niroo Research Institute}, Tehran, February 2020, and \textbf{MCI and the Iran Telecom Exhibition}, 2018 to 2019. Invited industry lectures on SDN, NFV, network slicing and 5G core migration.
  \item \textbf{Graduate mentoring}, 2024 to present. Advised more than fourteen international applicants on doctoral applications, resulting in nine successful placements.
\end{itemize}
\end{rSection}

%-----------------------------------------------------------------------
\begin{rSection}{Technical Skills}
\textbf{Networks:} 5G and 6G systems, network slicing, O-RAN, SDN and NFV, multi-access edge computing, 3GPP and ETSI standards, Wireshark, TRex, GNS3, ONOS.\\
\textbf{Machine learning:} deep reinforcement learning (Stable-Baselines3, Ray RLlib, Gymnasium), PyTorch, scikit-learn, convex and integer optimisation (CVX, Gurobi), supervised learning.\\
\textbf{Software and tools:} Python, SQL, MATLAB, C and C++, Linux, Git, Docker, Grafana, \LaTeX.
\end{rSection}

\end{document}
